% =====================================================================
% DRAFT SECTION FOR PAPER 47 -- NOT INTEGRATED INTO main.tex
% Editor draft, 2026-10-13. Owner decides.
% Sources: R186_024 (Sec. 1.2-1.4), R186_025 (Lemmas 1.1-1.4, 2.1, Thms A, B,
% Cor. (L1), Sec. 3.2), R186_100 (Thm 3.1 in form-domain version), hostile
% reviews R186_104, R186_105 (what enters and with which label).
% Every statement was checked against the paper's own definitions:
% eq:symbol, eq:form (form domain), eq:W, P_T, thm:A/eq:Aform, lem:K/eq:m,
% thm:S(iii), lem:mbounds(i)(ii)(iii), def:PD, rem:Kprim, subsec:arch
% (A_inf, Pi_T), the paragraph before lem:contA (H_T, compact resolvent),
% lem:contA, thm:margin. See CHECK_TABLE.md, part III.
% Macros are those of main.tex. Lines "% EDITOR FLAG" must be resolved
% before integration. Compiles standalone (the preamble below is skipped
% when the file is \input into main.tex, where \PT is already defined).
% =====================================================================
\ifdefined\PT\else
\documentclass[11pt]{amsart}
\usepackage{amsmath,amssymb,amsthm,mathtools,mathrsfs}
\usepackage[margin=1.08in]{geometry}
\usepackage{booktabs}
\usepackage[hypertexnames=false]{hyperref}
\hypersetup{hidelinks}
\newtheorem{theorem}{Theorem}[section]
\newtheorem{proposition}[theorem]{Proposition}
\newtheorem{lemma}[theorem]{Lemma}
\newtheorem{corollary}[theorem]{Corollary}
\theoremstyle{definition}
\newtheorem{definition}[theorem]{Definition}
\theoremstyle{remark}
\newtheorem{remark}[theorem]{Remark}
\newcommand{\C}{\mathbb C}
\newcommand{\R}{\mathbb R}
\newcommand{\Rel}{\operatorname{Re}}
\newcommand{\supp}{\operatorname{supp}}
\newcommand{\Kp}{\mathcal K^{\mathrm{prim}}_{2T}}
\newcommand{\PT}{\mathcal P_T}
\newcommand{\IT}{I_T}
\newcommand{\lamP}{\lambda^{\mathrm{prim}}_T}
\newcommand{\lamA}{\lambda_{\min}(A_T|\PT)}
\newcommand{\Fh}{\widehat F}
\newcommand{\cF}{c_F}
\newcommand{\Lev}{\mathfrak E}
\newcommand{\NT}{\mathcal N_T}
\newcommand{\Aconst}{\mathsf A_T}
\newcommand{\Ainf}{\mathcal A_\infty}
\newcommand{\PiT}{\Pi_T}
\newcommand{\DraftStandalone}{}
\begin{document}
\title{Draft section for Paper 47: archimedean coercivity on Tate-orthogonal blocks}
\maketitle
\fi

\section{Archimedean coercivity on Tate-orthogonal blocks and the finite reduction}\label{sec:coer}

This section localizes the positivity question for $\mathfrak a_T$ on $\PT$ to an explicit
subspace of dimension $2k-2$, the \emph{piecewise Tate exponentials}, by proving that the
archimedean form is coercive, with explicit constants, on the complementary subspace of functions
that satisfy the two Tate conditions \emph{on each of $k$ consecutive blocks}. The coercivity is a
high-pass phenomenon ($m(\tau)\sim\log|\tau|$, Lemma~\ref{lem:mbounds}); the prime form is handled
by a norm bound. Nothing in this section bears on the sign of $\mathfrak a_T$: the finite statement
of Theorem~\ref{thm:coer-schur} inherits the margin of Theorem~\ref{thm:margin} in full
(Remark~\ref{rem:coer-honest}).

\emph{Notation.} All forms are on $L^2(\IT)$ with its metric. The form domain of \eqref{eq:form} is
$\mathcal D_T=\{F\in L^2(\IT):\int\log(2+|\tau|)|\Fh|^2<\infty\}$. By Theorem~\ref{thm:A},
\[
\mathfrak a_T=\Gamma+\PiT,\qquad
\Gamma(F):=\frac1{2\pi}\int_\R m(\tau)|\Fh(\tau)|^2d\tau=C_0\|F\|^2+\Lev(\cF),
\]
\[
\PiT(F)=-\sum_{n\in\NT}v_n\Rel\cF(\log n),
\]
where $\Gamma$ is the form $\Ainf$ of \S\ref{subsec:arch} without its polar term ($\Gamma=\Ainf$ on
$\PT$) and does not depend on $T$. For $F,G\in\mathcal D_T$ we write $\Gamma(F,G)$,
$\mathfrak a_T(F,G)$ for the polarized (sesquilinear) forms. For a window $\tau>0$ we write
$\lambda_{\mathrm A}(\tau):=\lambda_{\min}(A_\tau|\mathcal P_\tau)$ (so $\lambda_{\mathrm A}(T)=\lamA$);
it is the bottom of $\mathfrak a_\tau$ on the unit sphere of $\mathcal P_\tau\cap\mathcal D_\tau$, it is
attained (the operator $H_\tau$ of the paragraph preceding Lemma~\ref{lem:contA} has compact
resolvent), and it is continuous in $\tau$ (Lemma~\ref{lem:contA}). For $2\tau\le\log2$ the contact
set $\mathcal N_\tau$ is empty, so $\mathfrak a_\tau=\Gamma$ on $L^2(I_\tau)$ and
$\lambda_{\mathrm A}(\tau)$ is a purely archimedean quantity. We use throughout that
$|\cF(u)|\le\cF(0)=\|F\|^2$ and that $\cF(u)=0$ for $|u|\ge2T$.
% EDITOR FLAG (same as part I of the check table): main.tex l.2316 writes Pi_T without Re; all
% uses here are for real F, and for complex F the Re is the one of thm:A.

\subsection{Blocks, the Tate-orthogonal subspace and the reduced kernel}\label{subsec:coer-blocks}

\begin{definition}\label{def:coer-V}
Let $k\ge1$ and let $B_1,\dots,B_k$ be the $k$ consecutive open intervals of length $2\tau$,
$\tau:=T/k$, that partition $\IT$ up to endpoints, $B_j=(-T+2(j-1)\tau,\,-T+2j\tau)$. Put
\[
\mathcal V_k:=\Big\{F\in L^2(\IT):\ \int_{B_j}F(t)e^{t/2}dt=\int_{B_j}F(t)e^{-t/2}dt=0\ \ (1\le j\le k)\Big\},
\qquad E_k:=\PT\ominus\mathcal V_k ,
\]
and $e_j^\pm:=e^{\pm t/2}\mathbf 1_{B_j}$. (Any partition into intervals of lengths $\le\log2$ could be
used in Lemmas~\ref{lem:coer-decomp}--\ref{lem:coer-sector}; equal blocks are needed only from
Lemma~\ref{lem:coer-lowpass} on.)
\end{definition}

\begin{lemma}[the decomposition $\PT=\mathcal V_k\oplus E_k$]\label{lem:coer-decomp}
\emph{(i)} $\mathcal V_k$ is a closed subspace of $\PT$. \emph{(ii)} The orthogonal complement of
$\mathcal V_k$ in $L^2(\IT)$ is $\operatorname{span}\{e_j^\pm\}$, of dimension $2k$.
\emph{(iii)} $E_k=\operatorname{span}\{e_j^\pm\}\cap\PT=\{\sum_jc_j^+e_j^++c_j^-e_j^-:\ A=B=0\}$ has
dimension $2k-2$, and $\PT=\mathcal V_k\oplus E_k$ (orthogonal in $L^2(\IT)$).
\emph{(iv)} $E_k\subset\mathcal D_T$, and $\PT\cap\mathcal D_T=(\mathcal V_k\cap\mathcal D_T)\oplus E_k$.
\end{lemma}
\begin{proof}
(i) The $2k$ conditions are continuous linear functionals; summing them over $j$ gives $A(F)=B(F)=0$.
(ii) $\mathcal V_k$ is the joint kernel of the functionals $F\mapsto\langle F,e_j^\pm\rangle$, and the $2k$
functions $e_j^\pm$ are linearly independent (disjoint supports; $e^{t/2},e^{-t/2}$ independent on an
interval). (iii) $E_k=\mathcal V_k^\perp\cap\PT$ since $\mathcal V_k\subset\PT$. The two functionals
$A,B$ restricted to $\operatorname{span}\{e_j^\pm\}$ are independent: on the two-dimensional space
$\operatorname{span}\{e_1^+,e_1^-\}$ their matrix is
$\begin{psmallmatrix}\int_{B_1}e^t&|B_1|\\|B_1|&\int_{B_1}e^{-t}\end{psmallmatrix}$, of determinant
$\int_{B_1}e^t\int_{B_1}e^{-t}-|B_1|^2>0$ (Cauchy--Schwarz, strict since $e^t$ is not constant). Hence
$\dim E_k=2k-2$. If $F\in\PT$ and $F=v+w$ with $v\in\mathcal V_k$, $w\in\mathcal V_k^\perp$, then
$w=F-v\in\PT$, so $w\in E_k$. (iv) $\widehat{e_j^\pm}(s)=\big[e^{(\pm\frac12+is)t}/(\pm\frac12+is)\big]_{\partial B_j}$,
so $|\widehat{e_j^\pm}(s)|\le2e^{T/2}(\frac14+s^2)^{-1/2}$ and $\int\log(2+|s|)|\widehat{e_j^\pm}|^2<\infty$.
Since $\mathcal D_T$ is a linear space containing $E_k$, the $\mathcal V_k$-component of an $F\in\PT\cap\mathcal D_T$
lies in $\mathcal D_T$.
\end{proof}

\begin{lemma}[the reduced kernel]\label{lem:coer-K1}
Put $K_1(u):=K(u)-2e^{-u/2}=\dfrac{2e^{-5u/2}}{1-e^{-2u}}=2\sum_{j\ge1}e^{-(2j+\frac12)u}$ ($u>0$), with
$K$ as in Lemma~\ref{lem:K}. Then $0<K_1(u)\le u^{-1}$ for all $u>0$, and $K_1(u)=u^{-1}-\frac32+O(u)$ as
$u\downarrow0$. Let $F,G\in L^2(\IT)$ be real, $\supp F\subset\overline B$, $\supp G\subset\overline{B'}$,
with $B,B'$ disjoint intervals and $B'$ to the right of $B$. Then
\begin{equation}\label{eq:coer-cross}
\Gamma(F+G)=\Gamma(F)+\Gamma(G)-\iint F(x)G(y)K(y-x)\,dx\,dy ,
\end{equation}
and if $\int F(x)e^{x/2}dx=0$ or $\int G(y)e^{-y/2}dy=0$, the kernel $K$ in \eqref{eq:coer-cross} may be
replaced by $K_1$.
\end{lemma}
\begin{proof}
The three expressions for $K_1$ agree by Lemma~\ref{lem:K} ($K=2e^{-u/2}\sum_{j\ge0}e^{-2ju}$).
$K_1\le u^{-1}$ is $e^{5u/2}-e^{u/2}\ge2u$: the function $g(u)=e^{5u/2}-e^{u/2}-2u$ has $g(0)=g'(0)=0$ and
$g''=\frac{25}4e^{5u/2}-\frac14e^{u/2}>0$. The expansion follows from
$\frac1{2\sinh(u/2)}=u^{-1}-\frac u{24}+O(u^3)$, $\frac1{2\cosh(u/2)}=\frac12+O(u^2)$ and
$2e^{-u/2}=2-u+O(u^2)$. For \eqref{eq:coer-cross}, apply \eqref{eq:Aform} with $\NT=\emptyset$ to
$F+G$: $c_{F+G}(u)=\cF(u)+c_G(u)+\int F(t)G(t-u)dt+\int G(t)F(t-u)dt$, and for $u>0$ the first cross
integral vanishes ($t\in B$, $t-u\in B'$ is impossible), while the second equals
$\int\!\!\int F(x)G(y)\mathbf 1_{y-x=u}$; integrating against $K(u)du$ over $u>0$ gives
\eqref{eq:coer-cross} (the term $\|F+G\|^2=\|F\|^2+\|G\|^2$ has no cross part). Finally
$2e^{-(y-x)/2}=2e^{-y/2}e^{x/2}$, so $\iint F(x)G(y)\,2e^{-(y-x)/2}dxdy=2\int Fe^{x/2}\int Ge^{-y/2}$,
which vanishes under either Tate condition.
\end{proof}
% EDITOR FLAG (normalization): the cross term in (eq:coer-cross) carries the factor 1 for the
% UNORDERED pair {B,B'}, i.e. the bilinear form Gamma(F,G) = -(1/2) iint F G K. R186_104 row 2 writes
% "-2 iint", and the scripts r025_coupling_norm.py compute Q_tot = 2<F_B, K_1 F_B'>, i.e. twice the
% cross term. See Remark rem:coer-nogo and check table III, conflict C1.

\begin{lemma}[a block sector is a short window]\label{lem:coer-sector}
Let $B\subset\IT$ be an interval of length $2\tau$ with $2\tau\le\log2$. Then
\[
\inf\Big\{\mathfrak a_T(F):\ F\in\mathcal D_T,\ \supp F\subset\overline B,\ \int_BFe^{\pm t/2}=0,\ \|F\|=1\Big\}
=\lambda_{\mathrm A}(\tau),
\]
independently of $T$ and of the position of $B$; the infimum is attained and depends continuously on
$\tau$. In particular, for $F\in\mathcal V_k$ with $\tau=T/k$, $2\tau\le\log2$, and $F_j:=F\mathbf 1_{B_j}$,
\begin{equation}\label{eq:coer-diag}
\sum_{j=1}^k\Gamma(F_j)\ \ge\ \lambda_{\mathrm A}(\tau)\,\|F\|^2 .
\end{equation}
\end{lemma}
\begin{proof}
Let $B=I_\tau+a$. For $F$ supported in $\overline B$ and $F_a(t):=F(t+a)$, supported in $\overline{I_\tau}$,
one has $c_{F_a}=\cF$, hence $\Gamma(F_a)=\Gamma(F)$ by \eqref{eq:Aform}, and
$\int F_ae^{\pm t/2}=e^{\mp a/2}\int Fe^{\pm t/2}$, so the two Tate conditions on $B$ are those of
$\mathcal P_\tau$ for $F_a$. Since $\cF(\log n)=0$ for $\log n\ge2\tau$ and $\log2\ge2\tau$, also
$\PiT(F)=0$, so $\mathfrak a_T(F)=\Gamma(F)=\mathfrak a_\tau(F_a)$ and the infimum is
$\lambda_{\mathrm A}(\tau)$. Attainment and continuity are the facts recalled in the notation paragraph.
\eqref{eq:coer-diag} is the sum over blocks.
\end{proof}

\subsection{A lower bound for the short window}\label{subsec:coer-short}

\begin{lemma}\label{lem:coer-K}
$K(u)\ge u^{-1}$ for $0<u\le2$.
\end{lemma}
\begin{proof}
By Lemma~\ref{lem:K}, $K(u)-u^{-1}=\frac{u-2\sinh(u/2)}{2u\sinh(u/2)}+\frac1{2\cosh(u/2)}$. From
$\tanh x\le x$, $2\sinh(u/2)\le u\cosh(u/2)$, so $u-2\sinh(u/2)\ge-u(\cosh(u/2)-1)=-2u\sinh^2(u/4)$ and
the first term is $\ge-\frac{\sinh(u/4)}{2\cosh(u/4)}\ge-\frac u8$. Hence
$K(u)-u^{-1}\ge\frac1{2\cosh(u/2)}-\frac u8\ge0$ whenever $u\cosh(u/2)\le4$, which holds for $u\le2$
($2\cosh1=3.09$).
\end{proof}

\begin{lemma}[the tail]\label{lem:coer-tail}
For $\tau>0$, $\int_{2\tau}^\infty K(u)\,du\ge\log\frac2\tau+\frac\pi2-\tau$.
\end{lemma}
\begin{proof}
By Theorem~\ref{thm:S}(iii) the integral is $\log\coth\frac\tau2+\pi-2\arctan e^\tau$; $\coth x\ge x^{-1}$
gives the first term, and $\frac{d}{d\tau}(\pi-2\arctan e^\tau)=-\frac1{\cosh\tau}\ge-1$ with value
$\frac\pi2$ at $\tau=0$ gives the second.
\end{proof}

\begin{lemma}[the truncated logarithm]\label{lem:coer-w}
Put $w(s):=\int_0^2\frac{1-\cos sv}v\,dv=\log(2|s|)+\gamma-\operatorname{Ci}(2|s|)$ for $s\ne0$, $w(0)=0$.
Then $w$ is even, increasing in $|s|$, and for every $0<R\le\pi$ and all $s\in\R$,
$w(s)\ge w(R)\min\{1,s^2/R^2\}$.
\end{lemma}
\begin{proof}
$w'(s)=\int_0^2\sin(sv)\,dv=\frac{1-\cos2s}s\ge0$ for $s>0$. Next $w(s)/s^2=\int_0^2v\,\frac{1-\cos sv}{(sv)^2}dv$,
and $x\mapsto\frac{1-\cos x}{x^2}=\frac12\big(\frac{\sin(x/2)}{x/2}\big)^2$ is decreasing on $(0,2\pi)$;
for $0<s\le R\le\pi$ and $v\le2$ one has $sv\le2\pi$, so $w(s)/s^2\ge w(R)/R^2$. For $s\ge R$ use monotonicity.
\end{proof}

\begin{lemma}[a double Poincar\'e inequality]\label{lem:coer-poinc}
Let $h\in L^2(\R)$ with $\supp h\subset[-1,1]$ and $h''\in L^2(\R)$ (distributional derivative). Then
$\|h''\|^2\ge\pi^2\|h'\|^2\ge\frac{\pi^4}4\|h\|^2$.
\end{lemma}
\begin{proof}
$h,h''\in L^2(\R)$ give $h\in H^2(\R)\subset C^1(\R)$, and the support condition gives $h(\pm1)=h'(\pm1)=0$.
The second inequality is the Dirichlet Poincar\'e inequality for $h\in H^1_0(-1,1)$ (first eigenvalue
$(\pi/2)^2$). For the first, $g:=h'\in H^1_0(-1,1)$ with $\int_{-1}^1g=h(1)-h(-1)=0$. Expand $g=\sum_{n\ge1}c_n\phi_n$
in the orthonormal Dirichlet eigenbasis $\phi_n(x)=\sin\frac{n\pi(x+1)}2$ of $(-1,1)$, so that
$\|g\|^2=\sum c_n^2$ and $\|g'\|^2=\sum(\frac{n\pi}2)^2c_n^2$. Since $\int\phi_n=\frac4{n\pi}$ for odd $n$ and
$0$ for even $n$, the constraint reads $c_1=-\sum_{n\ge3\ \mathrm{odd}}c_n/n$. The even-indexed terms already
satisfy $(\frac{n\pi}2)^2\ge\pi^2$, and for the odd-indexed ones
\[
\sum_{n\ \mathrm{odd}}\Big(\frac{n^2\pi^2}4-\pi^2\Big)c_n^2=\frac{\pi^2}4\Big[\sum_{n\ge3\ \mathrm{odd}}(n^2-4)c_n^2-3c_1^2\Big]\ \ge\ 0,
\]
because by Cauchy--Schwarz $c_1^2\le S\sum_{n\ge3\ \mathrm{odd}}(n^2-4)c_n^2$ with
$S=\sum_{n\ge3\ \mathrm{odd}}\frac1{n^2(n^2-4)}=\frac14\big(\frac43-\frac{\pi^2}8\big)=0.0249<\frac13$
(partial fractions and $\sum_{n\ \mathrm{odd}}n^{-2}=\pi^2/8$). Hence $\|g'\|^2\ge\pi^2\|g\|^2$. (The constant
$\pi^2$ is attained by $g=\phi_2$; this is the parity argument of R186\_105 in Fourier form: the constraint
touches only the even part of $g$.)
\end{proof}

\begin{theorem}[the short window]\label{thm:coer-A}
Let $0<\tau\le1$ with $2\tau\le\log2$, and put $a:=\tau^2/4$. For every $R\in(0,\pi]$,
\begin{equation}\label{eq:coer-A}
\lambda_{\mathrm A}(\tau)\ \ge\ C_0+\log\frac2\tau+\frac\pi2-\tau+w(R)\Big[1-\frac{16(R^2+a)^3}{27\pi^4R^2}\Big].
\end{equation}
In particular, with $R=2$ ($w(2)=\log4+\gamma-\operatorname{Ci}(4)=2.1045$),
\[
\lambda_{\mathrm A}(\tau)\ \ge\ \log\frac1\tau-1.209-\tau\qquad(0<\tau\le0.1),
\]
so $\lambda_{\mathrm A}(\tau)\ge\log\frac1\tau-1.31$ on $(0,0.1]$.
\end{theorem}
\begin{proof}
Let $F\in\mathcal P_\tau\cap\mathcal D_\tau$, $\|F\|=1$. Since $\mathcal N_\tau=\emptyset$ and $\cF(u)=0$ for
$u\ge2\tau$, \eqref{eq:Aform} gives
\[
\mathfrak a_\tau(F)=C_0+\int_{2\tau}^\infty K+\int_0^{2\tau}\big(1-\Rel\cF(u)\big)K(u)\,du .
\]
The tail is bounded by Lemma~\ref{lem:coer-tail}. Put $G(y):=\tau^{1/2}F(\tau y)$, so that $\supp G\subset[-1,1]$,
$\|G\|=1$, $\cF(u)=c_G(u/\tau)$, and
$\int_0^{2\tau}(1-\Rel\cF)K=\int_0^2(1-\Rel c_G(v))\,\tau K(\tau v)\,dv\ge\int_0^2\frac{1-\Rel c_G(v)}v\,dv$,
using $1-\Rel c_G\ge0$ and Lemma~\ref{lem:coer-K} ($\tau v\le2\tau\le2$). By Plancherel,
$1-\Rel c_G(v)=\frac1{2\pi}\int(1-\cos sv)|\widehat G(s)|^2ds$, so by Tonelli the last integral equals
$\frac1{2\pi}\int w|\widehat G|^2$, and by Lemma~\ref{lem:coer-w}, for $0<R\le\pi$,
\[
\frac1{2\pi}\int w|\widehat G|^2\ \ge\ w(R)\Big[1-\frac1{2\pi}\int\Big(1-\frac{s^2}{R^2}\Big)^+|\widehat G(s)|^2ds\Big].
\]
The Tate conditions $\int Fe^{\pm t/2}=0$ read $\int Ge^{\pm\tau y/2}=0$, i.e.\ $\widehat G(\mp i\tau/2)=0$.
$\widehat G$ is entire of exponential type $\le1$ and in $L^2(\R)$, so $H(s):=\widehat G(s)/(s^2+a)$ is entire of
exponential type $\le1$ and in $L^2(\R)$ ($s^2+a\ge a>0$ on $\R$); by the Paley--Wiener theorem $H=\widehat h$ with
$h\in L^2(\R)$, $\supp h\subset[-1,1]$, and $G=-h''+ah$, so $h''\in L^2$. Then
$1=\|G\|^2=\frac1{2\pi}\int(s^2+a)^2|H|^2=\|h''\|^2+2a\|h'\|^2+a^2\|h\|^2\ge\|h''\|^2\ge\frac{\pi^4}4\|h\|^2$
(Lemma~\ref{lem:coer-poinc}). On the other hand $(1-\frac{s^2}{R^2})^+(s^2+a)^2\le M_R:=\frac{4(R^2+a)^3}{27R^2}$
(the maximum of $(1-x/R^2)(x+a)^2$ on $0\le x\le R^2$ is at $x=(2R^2-a)/3$), so
$\frac1{2\pi}\int(1-s^2/R^2)^+|\widehat G|^2\le M_R\|h\|^2\le\frac{4M_R}{\pi^4}$. Collecting terms gives
\eqref{eq:coer-A}. For $R=2$ and $a\le\frac1{400}$ the bracket is $\ge0.9025$ and $w(2)\cdot0.9025\ge1.899$;
with $C_0+\log2+\frac\pi2=-3.1082$ this gives $\log\frac1\tau-1.209-\tau$.
\end{proof}

\begin{remark}\label{rem:coer-A}
(1) The only numerical input is the value $w(2)$ (a tabulated elementary quantity, $\operatorname{Ci}(4)=-0.1410$).
Optimizing $R$ in \eqref{eq:coer-A} gives $R=2.165$ and $\log\frac1\tau-1.181-\tau$ (numerical
supremum). (2) The true value is $\lambda_{\mathrm A}(\tau)=\log\frac1\tau-0.510\pm0.001$ for $\tau\in[0.02,0.1]$
(Galerkin, R186\_025 and R186\_104; NUMERICAL, upper bounds). The loss of the proof is in
$\int_{|s|<R}|H|^2\le\|h\|^2$ and in the constant $\pi^4/4$ (the clamped-beam constant $31.285$ would
improve it; not used). (3) Theorem~\ref{thm:coer-A} is a complement to Lemma~\ref{lem:contA} at the
left end: $\lambda_{\mathrm A}(\tau)\to+\infty$ as $\tau\to0$ at the rate $\log\frac1\tau$, consistently with
Corollary~\ref{cor:mono} ($\lambda_{\mathrm A}(\tau)\le C_0+\varphi((3\pi/2\tau)^2)=\log\frac1\tau+O(1)$).
(4) \emph{Extension to the primitive class (ARGUED, not proved).} For real $F$, Remark~\ref{rem:Kprim}
describes $\Kp$ by $F\in\PT$, or $A=A'=0$, or $B=B'=0$; in the last two cases $\widehat G$ has a double zero at
one of $\mp i\tau/2$ and $|\widehat G(s)|^2=(s^2+a)^2|H(s)|^2$ on $\R$ with the same $H$-structure, so the proof
goes through for those $F$. The elements of $\Kp$ coming from complex $F$ (Proposition~\ref{prop:weak}) are
not covered; the statement $\lambda^{\mathrm{prim}}_\tau\ge\log\frac1\tau-1.31$ is therefore not asserted.
\end{remark}

\subsection{Coercivity of the archimedean form on $\mathcal V_k$}\label{subsec:coer-B}

For $R>0$ let $P_R$ be the orthogonal projection of $L^2(\R)$ onto the band $\{|s|\le R\}$,
$\widehat{P_RF}=\Fh\mathbf 1_{|s|\le R}$.

\begin{lemma}[low-pass leakage of $\mathcal V_k$]\label{lem:coer-lowpass}
Let $F\in\mathcal V_k$ with equal blocks of half-length $\tau$, and $R>0$. Then
\[
\|P_RF\|^2\le\delta(R\tau)\,\|F\|^2,\qquad \delta(\theta):=\frac{\cosh^2\tau}{6}\Big(\theta^2+\frac{\tau^2}4\Big)^2 .
\]
\end{lemma}
\begin{proof}
$G:=P_RF$ is entire. Fix a block $B$ with centre $c$ and put, on $B$,
$u(x):=G(c)\cosh\frac{x-c}2+2G'(c)\sinh\frac{x-c}2\in\operatorname{span}\{e^{x/2},e^{-x/2}\}$, so that
$u(c)=G(c)$, $u'(c)=G'(c)$, $u''=u/4$. Then $r:=G-u$ satisfies $r(c)=r'(c)=0$ and
$r''-\frac r4=G''-\frac G4=:H$, whence $r(x)=\int_c^x2\sinh\frac{x-t}2\,H(t)\,dt$. By Cauchy--Schwarz and
$2\sinh\frac z2\le z\cosh\frac z2\le z\cosh\tau$ ($0\le z\le\tau$),
$|r(x)|^2\le\cosh^2\tau\,\frac{|x-c|^3}3\,\|H\|_{L^2(B)}^2$, so $\int_Br^2\le\cosh^2\tau\,\frac{\tau^4}6\|H\|_{L^2(B)}^2$.
Since $F_B:=F\mathbf 1_B$ is orthogonal to $e^{\pm x/2}\mathbf 1_B$, $\langle F_B,G\rangle=\langle F_B,r\rangle\le\|F_B\|\,\|r\|_{L^2(B)}$.
Summing over the blocks and using Cauchy--Schwarz once more,
\begin{align*}
\|G\|^2=\langle F,G\rangle&\le\|F\|\Big(\sum_B\|r\|_{L^2(B)}^2\Big)^{1/2}\le\|F\|\cosh\tau\,\frac{\tau^2}{\sqrt6}\,\|G''-\tfrac14G\|_{L^2(\R)}\\
&\le\|F\|\cosh\tau\,\frac{\tau^2}{\sqrt6}\Big(R^2+\frac14\Big)\|G\|,
\end{align*}
the first equality because $P_R$ is an orthogonal projection, the last inequality because
$\widehat G$ is supported in $|s|\le R$. Squaring gives the claim, as $\tau^4(R^2+\frac14)^2=((R\tau)^2+\frac{\tau^2}4)^2$.
\end{proof}

\begin{theorem}[coercivity on $\mathcal V_k$]\label{thm:coer-B}
Let $k\ge1$, $\tau=T/k\le1$, $\kappa:=\cosh^2\tau$, $a:=\tau^2/4$, $\delta$ as in
Lemma~\ref{lem:coer-lowpass}, $\theta_*:=(\sqrt{6/\kappa}-a)^{1/2}$ (the root of $\delta(\theta_*)=1$), and
$\ell(\rho):=\log\frac\rho{2\pi}-\frac1\rho-\frac1{2\rho^2}$ (the lower bound of Lemma~\ref{lem:mbounds}(ii)).
Then for every $F\in\mathcal V_k\cap\mathcal D_T$,
\begin{equation}\label{eq:coer-B}
\Gamma(F)\ \ge\ \Big(\log\frac1\tau-C_B(\tau)\Big)\|F\|^2,\qquad
C_B(\tau):=\log\frac1\tau-C_0\,\delta(\tau)-\frac{2\kappa}3\big[\Phi_\tau(\theta_*)-\Phi_\tau(\tau)\big],
\end{equation}
where $\Phi_\tau$ is the elementary primitive
\[
\Phi_\tau(\theta):=\Big(\frac{\theta^4}4+\frac{a\theta^2}2\Big)\big(\log\theta-\log(2\pi\tau)\big)-\frac{\theta^4}{16}-\frac{a\theta^2}4
-\tau\Big(\frac{\theta^3}3+a\theta\Big)-\frac{\tau^2\theta^2}4-\frac{a\tau^2}2\log\theta
\]
of $(\theta^3+a\theta)\,\ell(\theta/\tau)$. Moreover $C_B(\tau)\to\frac14-\log\frac{6^{1/4}}{2\pi}=1.6399$ as
$\tau\to0$, and $C_B(0.02)=1.6573$, $C_B(0.0619)=1.6957$, $C_B(0.1)=1.7331$ (evaluations of the closed form).
\end{theorem}
\begin{proof}
Let $\|F\|=1$ and $\mu:=\frac1{2\pi}|\Fh|^2ds$, a probability measure, so that $\Gamma(F)=\int m\,d\mu$ by
\eqref{eq:form} with $\NT=\emptyset$. By Lemma~\ref{lem:mbounds}(i) and \eqref{eq:m}, $m$ is even, real-analytic,
increasing on $[0,\infty)$ with $m(0)=C_0$; hence $m(s)=C_0+\int_0^{|s|}m'(\rho)d\rho$ with $m'\ge0$ and, by Tonelli,
\[
\int m\,d\mu=C_0+\int_0^\infty m'(\rho)\,\mu(|s|\ge\rho)\,d\rho\ \ge\ C_0+\int_0^{\rho_*}m'(\rho)\big(1-\delta(\rho\tau)\big)d\rho,
\qquad\rho_*:=\theta_*/\tau,
\]
since $\mu(|s|\ge\rho)=1-\|P_{\rho}F\|^2+\mu(|s|=\rho)\ge1-\delta(\rho\tau)$ by Lemma~\ref{lem:coer-lowpass} and the
integrand is nonnegative on $[0,\rho_*]$, where $\delta(\rho\tau)\le1$. Integrating by parts on $[0,\rho_*]$, with
vanishing boundary terms ($1-\delta(\theta_*)=0$ and $m(0)-C_0=0$),
\[
\int_0^{\rho_*}m'(\rho)(1-\delta(\rho\tau))d\rho=\int_0^{\theta_*}\big(m(\theta/\tau)-C_0\big)\delta'(\theta)\,d\theta,
\]
so that
\[
\int m\,d\mu\ \ge\ C_0\,\delta(0)+\int_0^{\theta_*}m(\theta/\tau)\,\delta'(\theta)\,d\theta ,
\]
using $\int_0^{\theta_*}\delta'=1-\delta(0)$. Here $\delta'(\theta)=\frac{2\kappa}3\theta(\theta^2+a)\ge0$. By
Lemma~\ref{lem:mbounds}(i),(ii), $m\ge\max\{C_0,\ell\}$ on $(0,\infty)$; we use $m\ge C_0$ on $[0,\tau]$ (i.e.\ $\rho\le1$)
and $m\ge\ell$ on $[\tau,\theta_*]$ (note $\ell(1)=-3.34>C_0$, and $\ell$ is increasing, so no further loss occurs
there). Hence
\[
\int m\,d\mu\ \ge\ C_0\,\delta(0)+C_0\big(\delta(\tau)-\delta(0)\big)+\int_\tau^{\theta_*}\ell(\theta/\tau)\,\delta'(\theta)\,d\theta
=C_0\,\delta(\tau)+\frac{2\kappa}3\big[\Phi_\tau(\theta_*)-\Phi_\tau(\tau)\big],
\]
which is \eqref{eq:coer-B}; the primitive is checked by differentiation. As $\tau\to0$: $\kappa\to1$, $a\to0$,
$\theta_*\to6^{1/4}$, $\delta(\tau)\to0$, $\frac{2\kappa}3\Phi_\tau(\tau)\to0$, and
$\frac{2\kappa}3\Phi_\tau(\theta_*)=\log\theta_*-\log(2\pi\tau)-\frac14+o(1)$ (the $a$- and $\tau$-terms are $o(1)$),
whence $C_B(\tau)\to\log2\pi-\log6^{1/4}+\frac14$. The three values are evaluations of the closed form (double precision).
\end{proof}

\begin{remark}\label{rem:coer-B}
(1) Using $m\ge\max\{C_0,\ell\}$ with the exact crossover instead of the split at $\rho=1$ changes $C_B$ by less than
$5\cdot10^{-5}$ on $(0,0.1]$. (2) The uniform statement ``$C_B(\tau)\le1.74$ on $(0,0.1]$'' holds on a grid of
twelve points where $C_B$ is increasing (R186\_105; NUMERICAL); monotonicity of $C_B$ is not proved and is not used:
Corollary~\ref{cor:coer-L1} only needs $C_B$ at the specific $\tau=T/k$. (3) The bound loses about $1.0$ against
the Galerkin value of $\inf\Gamma|_{\mathcal V_k}$ (e.g.\ $1.087$ versus $2.093$ at $T=\frac12\log5$, $k=13$,
NUMERICAL); the loss is in Lemma~\ref{lem:coer-lowpass}, whose constant is about $25$ times larger than needed
on random elements of $\mathcal V_k$ (R186\_105). (4) The theorem uses the whole multiplier $m$ and the Tate
conditions only through the geometry of $\mathcal V_k$ (Lemma~\ref{lem:coer-lowpass}); it does not separate the
block-diagonal part \eqref{eq:coer-diag} from the coupling \eqref{eq:coer-cross}. See Remark~\ref{rem:coer-nogo}
for why that separation is not viable.
\end{remark}

\begin{proposition}[the norm of the prime form]\label{prop:coer-prime}
\emph{(i)} $|\PiT(F)|\le\Aconst\|F\|^2$ for all $F\in L^2(\IT)$ and all $T$.
\emph{(ii)} For $\ell>0$ let $L(T,\ell):=\lceil2T/\ell\rceil$. Then $|\Rel\cF(\ell)|\le\cos\frac{\pi}{L(T,\ell)+1}\,\|F\|^2$,
hence
\[
\|\PiT\|\ \le\ \sum_{n\in\NT}v_n\cos\frac\pi{L(T,\log n)+1}
=\begin{cases}
\Aconst/2, & T\le\log2,\\[2pt]
\dfrac{v_2}{\sqrt2}+\dfrac{v_3+v_4}2=1.6740, & \log2<T\le\tfrac12\log5 .
\end{cases}
\]
Here $\frac{v_2}{\sqrt2}+\frac{v_3+v_4}2=\log2+\frac{\log3}{\sqrt3}+\frac{\log2}2$; in the first case the bound is
$\frac{v_2}2=0.4901$ for $T\le\frac12\log3$ and $\frac{v_2+v_3}2=1.1244$ for $\frac12\log3<T\le\log2$. The bounds hold on all of $L^2(\IT)$, in particular on $\mathcal V_k$.
\end{proposition}
\begin{proof}
(i) is $|\cF(\log n)|\le\|F\|^2$. (ii) Let $S_\ell$ be the shift $F\mapsto F(\cdot-\ell)\mathbf 1_{\IT}$ on
$L^2(\IT)$, so that $\Rel\cF(\ell)=\frac12\langle(S_\ell+S_\ell^*)F,F\rangle$. The map
$F\mapsto\big(F(-T+r+j\ell)\big)_{0\le j<n(r)}$, $r\in[0,\ell)$, $n(r):=\#\{j\ge0:-T+r+j\ell<T\}\le L(T,\ell)$, is unitary
from $L^2(\IT)$ onto $\int^\oplus_{[0,\ell)}\C^{n(r)}dr$, and it carries $S_\ell$ to the fibrewise shift $j\mapsto j+1$.
Hence $S_\ell+S_\ell^*$ acts on the fibre $\C^{n}$ as the adjacency matrix of the path with $n$ vertices, whose
eigenvalues are $2\cos\frac{j\pi}{n+1}$, $1\le j\le n$; its norm is $2\cos\frac\pi{n+1}\le2\cos\frac\pi{L+1}$. For
$T\le\log2$ one has $2T/\log n\le2$ for all $n\ge2$, so $L=2$ and $\cos\frac\pi3=\frac12$. For $\log2<T\le\frac12\log5$:
$\NT=\{2,3,4\}$, $L(T,\log2)=3$ since $2<2T/\log2\le\log5/\log2=2.32$, and $L(T,\log3)=L(T,\log4)=2$ since
$2T\le\log5<2\log3$; $\cos\frac\pi4=\frac1{\sqrt2}$. The values use $v_2=\sqrt2\log2$, $v_3=2\log3/\sqrt3$, $v_4=\log2$.
\end{proof}
% EDITOR FLAG: for (1/2)log2 < T < (1/2)log3 the bound A_T/2 is the exact norm (thm:inertia-B2 of the
% draft section sec:inertia, same fibre decomposition). R186_024 Sec. 3(c) claimed A_T/2 for all T; this is
% false for T > log 2 (orbit of log 2 of length 3), as found in R186_104 3b; the proposition is the corrected
% statement.

\begin{corollary}[coercivity of the truncated form on $\mathcal V_k$]\label{cor:coer-L1}
Let $k\ge1$ with $\tau=T/k\le1$, and put $c(T,k):=\inf\{\mathfrak a_T(F):F\in\mathcal V_k\cap\mathcal D_T,\ \|F\|=1\}$. Then
\begin{equation}\label{eq:coer-L1}
c(T,k)\ \ge\ \log\frac kT-C_B\Big(\frac Tk\Big)-\|\PiT\| ,
\end{equation}
with $\|\PiT\|$ bounded by Proposition~\ref{prop:coer-prime}. In particular (evaluations of the closed form):
\begin{center}
\begin{tabular}{@{}lllll@{}}
\toprule
$T$ & $k$ & $\dim E_k$ & bound for $\|\PiT\|$ & lower bound for $c(T,k)$\\
\midrule
$\frac12\log3$ & $8$ & $14$ & $0.4901$ & $0.486$\\
$\log2$ & $12$ & $22$ & $1.1244$ & $0.035$\\
$0.69$ & $14$ & $26$ & $1.1244$ & $0.202$\\
$\frac12\log5$ & $25$ & $48$ & $1.6740$ & $0.094$\\
$\frac12\log5$ & $26$ & $50$ & $1.6740$ & $0.134$\\
$\frac12\log5$ & $30$ & $58$ & $1.6740$ & $0.281$\\
$\frac12\log5$ & $85$ & $168$ & $\Aconst=2.9420$ & $0.070$\\
\bottomrule
\end{tabular}
\end{center}
Moreover, for every $T>0$ there is $k_0(T)$ such that $c(T,k)>0$ for all $k\ge k_0(T)$; indeed the right-hand side
of \eqref{eq:coer-L1} tends to $+\infty$ with $k$, because $C_B(T/k)\to1.6399$ (Theorem~\ref{thm:coer-B}).
\end{corollary}
\begin{proof}
$\mathfrak a_T=\Gamma+\PiT$ on $\mathcal D_T$, Theorem~\ref{thm:coer-B} and Proposition~\ref{prop:coer-prime}.
\end{proof}

\begin{remark}\label{rem:coer-size}
The size of $k_0(T)$: by \eqref{eq:coer-L1} and Proposition~\ref{prop:coer-prime}(i), $c(T,k)>0$ as soon as
$\log\frac kT>C_B(T/k)+\Aconst$, so $k_0(T)\le Te^{\Aconst+C_B(T/k_0)}$, where $C_B(T/k_0)$ is within $0.1$ of
its limit once $T/k_0\le0.1$ (NUMERICAL; the limit itself is proved). Since $\Aconst\asymp e^T$, the dimension
$2k-2$ of $E_k$ grows doubly exponentially in $T$; on the declared range $T\le\frac12\log5$ it stays below $60$.
The Galerkin values of $c(T,k)$ are larger than the bounds by about $1.0$ (Remark~\ref{rem:coer-B}(3)):
e.g.\ $c(\frac12\log5,13)=1.241$ (NUMERICAL, converged in the Galerkin degree), where \eqref{eq:coer-L1} gives
$-0.59$.
\end{remark}

\subsection{The finite reduction}\label{subsec:coer-schur}

\begin{theorem}[reduction to the piecewise Tate exponentials]\label{thm:coer-schur}
Fix $T>0$ and $k\ge1$ with $c:=c(T,k)>0$ (Corollary~\ref{cor:coer-L1}: every $k\ge k_0(T)$). Put
$\mathcal V:=\mathcal V_k\cap\mathcal D_T$ and $M:=\Aconst-C_0$.
\emph{(i)} $\mathfrak a_T$ is an inner product on $\mathcal V$ for which $\mathcal V$ is complete, and for every
$e\in E_k$ the functional $v\mapsto\mathfrak a_T(e,v)$ is bounded on $(\mathcal V,\mathfrak a_T)$; hence there is a unique
$Ke\in\mathcal V$ with $\mathfrak a_T(Ke,v)=\mathfrak a_T(e,v)$ for all $v\in\mathcal V$. The map $K:E_k\to\mathcal V$ is linear,
and $\|K\|:=\sup\{\|Ke\|:\|e\|=1\}<\infty$.
\emph{(ii)} For $v\in\mathcal V$ and $e\in E_k$,
\begin{equation}\label{eq:coer-schur}
\begin{split}
\mathfrak a_T(v+e)&=\mathfrak a_T(v+Ke)+\mathfrak s_{T,k}(e),\\
\mathfrak s_{T,k}(e)&:=\mathfrak a_T(e)-\mathfrak a_T(Ke)=\mathfrak a_T(e)-\sup_{u\in\mathcal V}\big[2\Rel\mathfrak a_T(e,u)-\mathfrak a_T(u)\big].
\end{split}
\end{equation}
Let $S_{T,k}$ be the matrix of the real quadratic form $\mathfrak s_{T,k}$ in an orthonormal basis of $E_k$ (a
real symmetric $(2k-2)\times(2k-2)$ matrix) and $\lambda_S$ its least eigenvalue. Then
\begin{equation}\label{eq:coer-sandwich}
\frac{\min\{c,\lambda_S\}}{(1+\|K\|)^2}\ \le\ \lamA\ \le\ \min\{c,\lambda_S\}.
\end{equation}
\emph{(iii)} Consequently $\lamA>0$ if and only if $S_{T,k}\succ0$, and $\lamA\ge0$ if and only if $S_{T,k}\succeq0$.
\end{theorem}
\begin{proof}
(i) By Lemma~\ref{lem:mbounds}, $\sigma_T\ge C_0-\Aconst$, so $\mathfrak a_T+M\|\cdot\|^2\ge0$ on $\mathcal D_T$, and
by Lemma~\ref{lem:mbounds}(ii) and $m\ge C_0$ there are constants $c_1,c_2>0$ (depending on $T$) with
$c_1\int\log(2+|s|)|\Fh|^2\le\mathfrak a_T(F)+(M+1)\|F\|^2\le c_2\int\log(2+|s|)|\Fh|^2$ on $\mathcal D_T$. On $\mathcal V$,
$\mathfrak a_T\ge c\|\cdot\|^2$ gives $\mathfrak a_T(v)\ge\frac c{c+M+1}\big(\mathfrak a_T(v)+(M+1)\|v\|^2\big)$, so
$\mathfrak a_T$ is equivalent on $\mathcal V$ to the weighted norm $\int\log(2+|s|)|\widehat v|^2$, for which
$\mathcal D_T$ is complete and $\mathcal V_k\cap\mathcal D_T$ is closed ($\mathcal V_k$ is $L^2$-closed and the weighted norm
dominates the $L^2$ norm). For $e\in E_k\subset\mathcal D_T$ (Lemma~\ref{lem:coer-decomp}(iv)), the Cauchy--Schwarz
inequality for the nonnegative form $\mathfrak a_T+M\langle\cdot,\cdot\rangle$ gives
$|\mathfrak a_T(e,v)|\le(\mathfrak a_T(e)+M\|e\|^2)^{1/2}(\mathfrak a_T(v)+M\|v\|^2)^{1/2}+M\|e\|\|v\|\le C_e\,\mathfrak a_T(v)^{1/2}$,
so $Ke$ exists by the Riesz representation theorem; $K$ is linear, and $\|K\|<\infty$ since $\dim E_k<\infty$.
(ii) $\mathfrak a_T(v+Ke)=\mathfrak a_T(v)+2\Rel\mathfrak a_T(Ke,v)+\mathfrak a_T(Ke)=\mathfrak a_T(v)+2\Rel\mathfrak a_T(e,v)+\mathfrak a_T(Ke)$,
which rearranges to \eqref{eq:coer-schur}; the variational expression is $\mathfrak a_T(Ke)=\sup_u[2\Rel\mathfrak a_T(Ke,u)-\mathfrak a_T(u)]$
(maximum at $u=Ke$). For \eqref{eq:coer-sandwich}: every $F\in\PT\cap\mathcal D_T$ is $v+e$ (Lemma~\ref{lem:coer-decomp}(iv)),
$\|F\|^2=\|v\|^2+\|e\|^2$, and $v\mapsto v+Ke$ is a bijection of $\mathcal V$. Taking $e=0$ gives $\lamA\le c$; taking
$v=-Ke$ gives $\mathfrak a_T(F)=\mathfrak s_{T,k}(e)$ with $\|F\|^2\ge\|e\|^2$, so $\lamA\le\lambda_S$. Conversely, if
$\lambda_S>0$, then $\mathfrak a_T(v+e)\ge c\|v+Ke\|^2+\lambda_S\|e\|^2$, and
$\|v\|^2+\|e\|^2\le(\|v+Ke\|+\|K\|\|e\|)^2+\|e\|^2\le(1+\|K\|)^2(\|v+Ke\|^2+\|e\|^2)$, which gives the lower bound; if
$\lambda_S\le0$ the lower bound is implied by $\lamA\le\lambda_S$. (iii) follows from \eqref{eq:coer-sandwich}
and, for the nonstrict case, from $\mathfrak a_T(v+e)\ge\mathfrak s_{T,k}(e)$ together with $v=-Ke$.
\end{proof}

\begin{remark}[what the reduction is and is not]\label{rem:coer-honest}
(1) Theorem~\ref{thm:coer-schur} is the Schur complement of $\mathfrak a_T$ with respect to $\PT=\mathcal V_k\oplus E_k$,
in the form-domain setting required because $\mathfrak a_T$ is unbounded on $L^2$ (R186\_100). Its content for the
positivity question is a \emph{localization}: the sign of $\lamA$ is the sign of an explicit real symmetric matrix of
size $2k-2$, with $k=25$ at $T=\frac12\log5$, whose index set is the piecewise Tate exponentials
$e^{\pm t/2}\mathbf 1_{B_j}$. The Galerkin near-minimizers of $\mathfrak a_T$ on $\PT$ have $99\%$ or more of their mass in
$E_k$ for $k\ge5$ (NUMERICAL, R186\_024/104); this is approximation theory (piecewise Tate exponentials approximate
smooth functions to $O(\tau^2)$), not arithmetic structure.
(2) \emph{No gain in precision.} By \eqref{eq:coer-sandwich}, $\lamA\le\lambda_S\le(1+\|K\|)^2\lamA$ whenever
$\lambda_S\le c$; so $S_{T,k}$ inherits the decay of Theorem~\ref{thm:margin} and the smallness already present at
the declared range: at $T=\frac12\log5$, $\lambda_S\approx\lamA\approx3\cdot10^{-12}$ with entries of $S_{T,k}$ of size
$O(1)$ (NUMERICAL, Galerkin; the certified value is Theorem~\ref{thm:log5A}). The entries of $S_{T,k}$ are not in
closed form: by \eqref{eq:coer-schur} the correction $\mathfrak a_T(Ke)$ is a supremum over $\mathcal V$, so Galerkin in
$\mathcal V$ gives lower bounds for the correction, i.e.\ upper bounds for $S_{T,k}$; a certified lower bound for $S_{T,k}$
needs a near-maximizer $u_N$ together with the dual estimate $\|\mathfrak a_T(e,\cdot)-\mathfrak a_T(u_N,\cdot)\|^2_{\mathcal V^*}/c$,
which is legitimate precisely because $c>0$ is proved. A certification by interval arithmetic along these lines is
realistic for $T\le0.6$, where $\lamA\ge10^{-5}$ (NUMERICAL, R186\_104/105), and is not attempted here.
(3) The decomposition with a \emph{fixed} number of blocks (the orbit blocks of the first contact, $k=3$, $\dim E=4$)
is not stable in $T$: the Galerkin value of $c(T,3)$ falls to $0.013$ at $T=0.69$ (NUMERICAL, R186\_024/104); only
$k\asymp T/\tau_0$ equal blocks give a uniform $c$ on a bounded range, by Corollary~\ref{cor:coer-L1}.
\end{remark}

\begin{remark}[the route by operator norms is not viable]\label{rem:coer-nogo}
By Lemmas~\ref{lem:coer-sector} and \ref{lem:coer-K1}, for $F\in\mathcal V_k$,
$\Gamma(F)=\sum_j\Gamma(F_j)+\mathfrak q(F)$ with $\sum_j\Gamma(F_j)\ge\lambda_{\mathrm A}(\tau)\|F\|^2$ and the coupling form
$\mathfrak q(F)=-\sum_{i<j}\iint F_i(x)F_j(y)K_1(|x-y|)\,dxdy$, so that
$c(T,k)\ge\lambda_{\mathrm A}(\tau)-\sup_{\mathcal V_k}|\mathfrak q|/\|\cdot\|^2-\|\PiT\|$. This splitting is sharp in neither
term and cannot replace Theorem~\ref{thm:coer-B}:
\begin{itemize}
\item[(a)] For two adjacent blocks the bilinear form $(F_1,F_2)\mapsto\iint F_1F_2K_1(|x-y|)$ on Tate-orthogonal
functions has norm exactly $\pi$: $\le\pi$ by $0<K_1(u)\le u^{-1}$ and Hilbert's inequality
$\iint\frac{f(x)g(y)}{x+y}\le\pi\|f\|\|g\|$ on $(0,\infty)$ (CITED, \cite[Thm.~316]{HLP}); $\ge\pi$ by the Weyl sequence
$f_M(x')=x'^{-1/2}\mathbf 1_{[2\tau e^{-M},2\tau]}$ in the distance $x'$ to the common endpoint, placed on both blocks and
projected onto the orthogonal complement of $\{e^{\pm x/2}\mathbf 1_B\}$: in logarithmic variables
$\iint f_Mf_M/(x'+y')=\int_{[-M,0]^2}\frac{du\,dv}{2\cosh\frac{u-v}2}=\pi M-O(1)$, $\|f_M\|^2=M$, the bounded part
$K_1-u^{-1}=-\frac32+O(u)$ contributes $O(\|f_M\|_1^2)=O(\tau)$, and the Tate projection changes $f_M$ by a vector of norm
$O(1)$ uniformly in $M$ (moments $O(\sqrt\tau)$, Gram matrix $\asymp\tau$), which changes the pairing by $O(\sqrt M)$
(R186\_025 \S3.2, verified in R186\_105). Hence $\sup_{\mathcal V_2}|\mathfrak q|/\|\cdot\|^2=\frac\pi2$ exactly, and
$\sup_{\mathcal V_k}|\mathfrak q|/\|\cdot\|^2\ge\frac\pi2$ for every $k\ge2$ and every $\tau$; the adjacent part of
$\mathfrak q$ is $\le\pi\cos\frac\pi{k+1}\|F\|^2$ by the path bound of Proposition~\ref{prop:coer-prime}, and the
non-adjacent part is not estimated here.
\item[(b)] The value $\frac\pi2$ is a property of the essential spectrum of the Carleman operator $1/(x+y)$ (symbol
$\pi/\cosh\pi\sigma$ under the Mellin transform): no condition of finite codimension, such as the Tate conditions,
lowers it. Galerkin values of this norm ($0.74$ at degree $10$, R186\_024) are what polynomials of low degree see and
grow with the degree without saturating (NUMERICAL, R186\_025/104).
\item[(c)] With \eqref{eq:coer-A}, the split bound is positive at $T=\frac12\log5$ only for
$\log\frac1\tau>1.21+\tau+\frac\pi2+1.674$, i.e.\ $k\ge70$ ($\dim E_k\ge138$), against $k=25$ from
Theorem~\ref{thm:coer-B}; and the true coupling cost at the Galerkin minimizers of $\Gamma|_{\mathcal V_k}$ is only
$0.17$--$0.18$ (NUMERICAL, R186\_025): the extremals of the coupling norm are not the minimizers of the form. The
relative bound $\mathfrak q\le\theta\sum_j\Gamma(F_j)$ holds numerically with $\theta\approx0.22$ (NUMERICAL, not proved)
and is not needed: Theorem~\ref{thm:coer-B} is the relative bound with $\theta=1$ and an explicit constant.
\end{itemize}
\end{remark}
% EDITOR FLAG (constant pi versus pi/2): R186_024/025/104/105 state "||K_1||_{V_k} = pi". With the normalization
% of (eq:coer-cross) this pi is the norm of the BILINEAR form with kernel K_1 (factor 1), whose contribution to
% Gamma(F) is the quadratic form q with sup |q|/||F||^2 = pi/2 on the pair. The scripts (r025_coupling_norm.py,
% "Q_tot = 2 <f_B, H f_B'>") compute twice q. The conclusion (NO-GO) is unchanged; the quoted constant in the
% memos should read pi/2 for the quadratic form and pi for the bilinear one. Bib entry needed: HLP = Hardy,
% Littlewood, Polya, Inequalities, 2nd ed., Cambridge 1952 (Thm. 316 is Hilbert's inequality).

\begin{remark}[labels]\label{rem:coer-labels}
Lemmas~\ref{lem:coer-decomp}--\ref{lem:coer-poinc}, Theorems~\ref{thm:coer-A}, \ref{thm:coer-B}, \ref{thm:coer-schur},
Proposition~\ref{prop:coer-prime} and Corollary~\ref{cor:coer-L1} are proved, with the numerical inputs
$w(2)=2.1045$, the closed-form values of $C_B$, and the arithmetic of the table (double precision evaluations of
elementary expressions). The following are numerical and not proved: $\lambda_{\mathrm A}(\tau)=\log\frac1\tau-0.510$;
the uniform bound $C_B\le1.74$ on $(0,0.1]$; the Galerkin values of $c(T,k)$, $\lambda_S$, $\|K\|$, of the mass of the
near-minimizer in $E_k$, of the coupling cost at the minimizers and of the relative constant $\theta\approx0.22$.
Remark~\ref{rem:coer-A}(4) is argued, not proved. The positivity of $S_{T,k}$ in the declared range beyond what
Theorem~\ref{thm:log5A} already certifies for $\lamA$ is open; the route of Remark~\ref{rem:coer-nogo} is closed.
\end{remark}

\ifdefined\DraftStandalone\end{document}\fi
